% r03_hyb_gain: proof of \Cref{rot:prop:BD}, part VII: the gain of every difference element at a non-zero gauge field (the lift of a constant field and the phases, the contour-line comparison, contour vertices, contour lines, double contour lines, the list); this completes part (b).

\paragraph{The difference elements at a non-zero gauge field.}
After the constant-background reduction of \Cref{rot:lem:CB}, applied to both flows on their infinite lattices with the torus lines as exact periodisations, every propagator that is compared between the flows is an $A=0$ line $G_r(\Lambda_F,0)=\sum_{j<r}S^F_j$, exactly the kernels of \Cref{rot:prop:fc,rot:lem:H,rot:lem:slice,rot:cor:unif}, with all vertices in the region $F_v$, and the background enters only through the explicit coarse fields $A_r(w_1)\cdots A_r(w_q)$ of \cite[(3.55)]{K1} (the same for both flows), the constant-background phases \cite[(3.46)]{K1}, the contour vertices \cite[(3.24)--(3.26)]{K1} with their contour lines \cite[(3.27)]{K1}, the external lifted fields $\tilde A^{(r),F}$, and the per-flow remainders, which are bounded by \Cref{rot:prop:BD}(a) at the regular fields of \Cref{rot:lem:CB}(R). The elements that are not $A=0$ lines are compared here.

\begin{lemma}[the lift of a constant field; the phases at matched vertices]\label{rot:hyb:lem:bph}
The lift \cite[(2.18)]{K1} of a constant coarse field $B$ to either flow is the constant fine field $\vr_rB$, $\vr_r\defeq a_r/(a_r+\lmu^2\delta_r^2)\le1$, the same number in both flows (\cite[(3.43)]{K1}). For matched $v'\in\Lambda_X$, $v''=\tau v'$ and the common root $y\in\Gamma_r$,
\[
\big|U(\vr_rB(\Gamma^X_{y,v'}))-U(\vr_rB(\Gamma^Y_{y,v''}))\big|\le e_r|B|\,c_\tau\eta'\le e\,p_r\lmu^{-1}c_\tau\,\delta_r^{-1/2}\eta'
\]
for $B$ in the small-field class.
\end{lemma}
\begin{proof}
$Q_F^*\mathbf 1_{\Gamma_r}=\mathbf 1_{\Lambda_F}$ for the frame-covariant photon average (\Cref{rot:lem:tile}(iii): $(\hat Q^*\psi)(x)=\sum_\alpha\psi_\alpha(y(x))f_\alpha$, which for the frame components $\psi_\alpha=f_\alpha\cdot B$ of a constant $B$ is $B$), and $(-\Delta+\lmu^2\delta_r^2)\mathbf 1=\lmu^2\delta_r^2\mathbf 1$, so $(M+a_rQ_F^*Q_F)\mathbf 1=(\lmu^2\delta_r^2+a_r)\mathbf 1$ and $a_rG^{F,\mathrm{ph}}_rQ_F^*B=\vr_rB$ for both flows. Both phases are $e^{\mathrm ie_r\vr_rB\cdot(v-y)}$ with $v=v'$, $v''$ (\Cref{rot:hyb:lem:alg}(iii)); $|v'-v''|\le c_\tau\eta'$, $\vr_r\le1$, $|B|\le p_r(\lmu\delta_r)^{-1}$ (\cite[(3.2)]{K1}), $e_r=e\delta_r^{1/2}$.
\end{proof}

\begin{lemma}[contour-line comparison]\label{rot:hyb:lem:CL}
Let $\cA$ be a $C^1$ vector field on a neighbourhood of the level-$r$ block $z+R^{-1}[-\frac12,\frac12]^3$ with $|\cA|\le P_0$, $|\nabla\cA|\le P_1$ there. Let $a^X_b\defeq\cA(x_b)\cdot e_\mu$ for the $\Lambda_X$-bond $b=(x_b,x_b+\eta'e_\mu)$ and $a^Y_{b''}\defeq\cA(x_{b''})\cdot f_\alpha$ for the $\Lambda_Y$-bond $b''=(x_{b''},x_{b''}+\eta'f_\alpha)$ (site sampling, King's identification of bond and site fields). Then for $x'\in B^X_r(z)$, $x''=\tau x'$,
\[
\big|a^X(\Gamma^{X,(r)}_{z,x'})-a^Y(\Gamma^{Y,(r)}_{z,x''})\big|\ \le\ 10\,(P_0+P_1)\,\eta' .
\]
The same holds for the partial contours $\Gamma^{(j+1)}_{x_{j+1},x}$ of \cite[(2.19)]{K1}, with $P_0,P_1$ taken on the level-$(j+1)$ block, the sliver area $\le1.5\eta'\ell_{j+1}+33\eta'^2$ in the flux term and the lengths $\le1.5\ell_{j+1}$ ($Y$), $\le2.1\ell_{j+1}$ ($X$) in the Riemann term: the bound is $\le(7.1P_0+(6.6\ell_{j+1}+66\eta')P_1)\eta'\le C(P_0+P_1\ell_{j+1})\eta'$.
\end{lemma}
\begin{proof}
\emph{Riemann step.} For a bond $b$ of length $\eta'$ sampled at its initial point, $|\eta'\,\cA(x_b)\cdot\hat b-\int_b\cA\cdot dl|\le\eta'\cdot\eta'P_1$. Summing over the $\le2.1/\eta'$ bonds of $\Gamma^X$ and the $\le1.5/\eta'$ bonds of $\Gamma^Y$: $|a^X(\Gamma^X)-\int_{\Gamma^X}\cA|+|a^Y(\Gamma^Y)-\int_{\Gamma^Y}\cA|\le3.6\,P_1\eta'$. \emph{Stokes step.} $\int_{\Gamma^X}\cA-\int_{\Gamma^Y}\cA=\oint_{\partial\Sigma}\cA-\int_{[x',x'']}\cA=\int_\Sigma\curl\cA\cdot dS-\int_{[x',x'']}\cA$ with the surface $\Sigma=\Sigma(z,x')$ of \Cref{rot:hyb:lem:len}(iii) (piecewise planar, so Stokes' theorem applies piece by piece and the interior edges cancel); $|\curl\cA|\le2P_1$, $\mathrm{area}(\Sigma)\le1.5\eta'+33\eta'^2\le3\eta'$ for $r\ge2$, and $|[x',x'']|\le c_\tau\eta'$. Total $\le(3.6+6)P_1\eta'+7.1P_0\eta'\le10(P_0+P_1)\eta'$. For partial contours the same decomposition applies to the sub-polygon up to level $j+1$ (\Cref{rot:hyb:lem:len}(i),(iii)).
\end{proof}

\begin{lemma}[contour vertices at the lifted background]\label{rot:hyb:lem:CV}
Let $A'^F=\vartheta\tilde A^{(r),F}+za_F$ be the reduced background of the flow $F$ (\Cref{rot:lem:CB}). Then for matched $(y,x')$, $(y,x'')$,
\[
\big|U(A'^X(\Gamma^{X,(r)}_{y,x'}))-U(A'^Y(\Gamma^{Y,(r)}_{y,x''}))\big|\ \le\ e_r\big|A'^X(\Gamma^X_{y,x'})-A'^Y(\Gamma^Y_{y,x''})\big|\ \le\ C\,e_r\,p_r^2\,\eta^{1/2},
\]
and hence the kernels of $F^X_{2,r}(A'^X,0)$ and of $F^Y_{2,r}(A'^Y,0)\circ\tau$ (\cite[(3.51)]{K1}, the contour vertices (3.24)--(3.26)) differ entrywise by $\le Ce_rp_r^2\eta^{1/2}$ against their own size $\le Ce_rp_r^2$: a relative gain $\eta^{1/2}$, and an absolute gain $\eta^{1/2}e\,p_r^2\delta_r^{1/2}$. The products (3.24)--(3.26)$\times$(3.26) at the same $y$ inherit the bound.
\end{lemma}
\begin{proof}
Let $\cA\defeq\vartheta\,E^S_{\mathrm{ph}}\tilde A_r+za^{\mathrm c}$, where $E^S_{\mathrm{ph}}\tilde A_r\defeq I_S[H^{Y,\mathrm{ph}}\tilde A_r]$ is the cubic-Hermite extension of the photon lift of $\tilde A_r$ along the flow $Y$ (\Cref{rot:hyb:def:ES}; componentwise in any frame) and $a^{\mathrm c}$ the smooth Cameron--Martin form (\Cref{rot:lem:CM}). $\cA$ is $C^1$ (\Cref{rot:hyb:lem:ESH}(iv)). By \Cref{rot:hyb:lem:ESH}(iii) and \Cref{rot:lem:CB}(R), $|E^S_{\mathrm{ph}}\tilde A_r(r)|\le\max_\sigma|\tilde A^{(r),Y}(y_\sigma)|\le Cp_r(|r-x_0|+\sqrt3\eta')$ and $|\nabla E^S_{\mathrm{ph}}\tilde A_r(r)|\le\frac{3\sqrt3}2\max_{\mathrm{edges}}|\partial\tilde A^{(r),Y}|\le Cp_r$; with \Cref{rot:lem:CM} for $a^{\mathrm c}$, $P_0,P_1\le Cp_r^2$ on the vertex region (of diameter $\le Cp_r$). Then $a^Y_{b''}=\cA(x_{b''})\cdot f_\alpha$ exactly (\Cref{rot:hyb:lem:ESB}(c); \Cref{rot:lem:CM} for $a_Y$), while
\[
|a^X_b-\cA(x_b)\cdot e_\mu|\le\vartheta\sum_z|H^{X,\mathrm{ph}}(x_b,z)-E^S_{\mathrm{ph}}(x_b,z)||\tilde A_r(z)|+|a_{X,\mu}(x_b)-a^{\mathrm c}\cdot e_\mu|\le C(C^S_EJ_1\eta)^{1/2}p_r^2+C\eta'^2
\]
by \Cref{rot:hyb:cor:DF}(a) at the $X$ site $x_b$ (the decaying form; $|\tilde A_r(z)|\le p_r(|z-x_0|+1)$, the contributions from far away as in \Cref{rot:lem:CB}(R)) and \Cref{rot:lem:CM}. Summing the lift mismatch with weight $\eta'$ over the $\le2.1/\eta'$ bonds of $\Gamma^X$ gives $\le C\eta^{1/2}p_r^2$, and \Cref{rot:hyb:lem:CL} gives $\le10(P_0+P_1)\eta'\le Cp_r^2\eta'$ for the rest. Multiply by $e_r$ for the phases ($|U(a)-U(a')|\le e_r|a-a'|$), and by $\eta'^3$ per summand of (3.51).
\end{proof}

\begin{lemma}[{the contour lines \cite[(3.27)]{K1}}]\label{rot:hyb:lem:CLN}
For $0\le j<r$ let $S^{F,\mathrm{ph}}_j$ be the stage-$j$ photon slice kernel (an $A=0$ kernel: the scalar slice $S^F_j$ times the identity, frame-free), and let $\cL^F_j(x;y)\defeq\sum_{b_1\in\Gamma^{F,(j+1)}_{x_{j+1},x}}\eta'\,S^F_j(x_{b_1},y)\,\hat b_1\in\R^3$ ($\hat b_1$ the direction of $b_1$, $y\in\Lambda_F$) be King's contour line (3.27) as a vector: its component along the direction of the photon bond $b=(y,y+\eta'\hat e)$ is $G_{(j)}(\Gamma^{(j+1)}_{x_{j+1},x},b)$ of \cite[(3.27)]{K1}. Lines are compared as vectors (\Cref{rot:lem:vert}(iii)). Then for matched $(x',y')$, $(x'',y'')=(\tau x',\tau y')$,
\[
\big|\cL^X_j(x';y')-\cL^Y_j(x'';y'')\big|\ \le\ C\Big(\frac{\eta'}{\ell_j}\Big)^{\theta}e^{-\delta\dist(\Gamma^{(j+1)},y')/\ell_j},\qquad\theta=\tfrac12,
\]
against the unweakened bound $Ce^{-\delta\dist/\ell_j}$ of each line: a slice-$j$ line weakened by the degree $\theta$ in the sense of \Cref{rot:thm:H1S}, with the prefactor $\eta'^\theta$. With one lattice derivative at the lattice end $y$ (the gradients of the two lines read as vectors in the common frame, \Cref{rot:hyb:lem:DS}) the same holds with one more factor $\ell_j^{-1}$ and, in place of $(\eta'/\ell_j)^{1/2}$, the gain $(\eta'/\ell_j)^{\theta_\partial}$ of \Cref{rot:hyb:lem:DS}, $\theta_\partial=\alpha/(2(2+\alpha))\ge\theta_1$: a slice-$j$ line with one derivative, weakened by the degree $\theta_\partial$.
\end{lemma}
\begin{proof}
Keep the second argument of the $Y$ kernel at the lattice point $y''$ throughout (no extension in the second argument is used) and write
\[
\cL^X_j(x';y')-\cL^Y_j(x'';y'')=\sum_{b_1\in\Gamma^X}\eta'\big[S^X_j(x_{b_1},y')-\tilde S^S_j(x_{b_1},y'')\big]\hat b_1
+\Big[\sum_{b_1\in\Gamma^X}\eta'\tilde S^S_j(x_{b_1},y'')\hat b_1-\sum_{b''_1\in\Gamma^Y}\eta'S^Y_j(x_{b''_1},y'')\hat b''_1\Big].
\]
\emph{First sum} ($2\le j<r$): \Cref{rot:hyb:cor:DF}(b) at the $X$ sites $x_{b_1}$, per bond $\le C(\eta'/\ell_j)^{1/2}\ell_j^{-1}e^{-\delta|x_{b_1}-y'|/\ell_j}$; with weight $\eta'$ over the $\le10.5\ell_j/\eta'$ bonds of the partial $X$ contour (\Cref{rot:hyb:lem:len}(i)): $\le C(\eta'/\ell_j)^{1/2}e^{-\delta\dist(\Gamma^X,y')/\ell_j}$. \emph{Bracket}: for a unit vector $u$, $u\cdot[\,\cdot\,]=a^X(\Gamma^X)-a^Y(\Gamma^Y)$ for the field $\cA_u(r)\defeq\tilde S^S_j(r,y'')\,u$, which is $C^1$ (\Cref{rot:hyb:lem:ESB}(c)), whose $Y$ samples are exactly $S^Y_j(x_{b''},y'')\,u\cdot f_\alpha$, and which has $P_0\le C\ell_j^{-1}e^{-\delta\dist(\Gamma,y'')/\ell_j}$, $P_1\le C\ell_j^{-2}e^{-\delta\dist(\Gamma,y'')/\ell_j}$ on the level-$(j+1)$ block (\Cref{rot:hyb:lem:ESB}(b); the block's diameter $\sqrt3\ell_{j+1}$ and $|y'-y''|\le c_\tau\eta'$ change only the constant). \Cref{rot:hyb:lem:CL} for partial contours gives $\le C(P_0+P_1\ell_{j+1})\eta'\le C(\eta'/\ell_j)e^{-\delta\dist/\ell_j}$ (flux, end segment and Riemann terms together: the kernel difference along the contour plus the flux of the slice kernel through the sliver, as \Cref{rot:cor:unif}(c) states); $u$ over an orthonormal basis gives the vector bound. The matching $y'\leftrightarrow y''$ of the second argument sits inside \Cref{rot:hyb:cor:DF}(b). \emph{The derivative version} ($2\le j<r$). With the lattice derivative on the lattice end, split the first sum at the matched $Y$ point: $\sum_{b_1\in\Gamma^X}\eta'\big[\nabla^X_yS^X_j(x_{b_1},y')-\nabla^Y_yS^Y_j(\tau x_{b_1},y'')\big]\hat b_1+\sum_{b_1\in\Gamma^X}\eta'\big[\nabla^Y_yS^Y_j(\tau x_{b_1},y'')-\nabla^Y_y\tilde S^S_j(x_{b_1},y'')\big]\hat b_1$ (gradients as vectors, the matrix $[\,\cdot\,]\hat b_1$ contracted on the line's index). The first sum is \Cref{rot:hyb:lem:DS} at the matched pairs $(x_{b_1},y')$, $(\tau x_{b_1},y'')$, per bond $\le C(\eta'/\ell_j)^{\theta_\partial}\ell_j^{-2}e^{-\delta|x_{b_1}-y'|/\ell_j}$, hence $\le C(\eta'/\ell_j)^{\theta_\partial}\ell_j^{-1}e^{-\delta\dist(\Gamma^X,y')/\ell_j}$ after the sum over the $\le10.5\ell_j/\eta'$ bonds with weight $\eta'$; the second sum is \Cref{rot:hyb:cor:DF}(c) per bond, $\le C(\eta'/\ell_j)^{1/2}\ell_j^{-1}e^{-\delta\dist/\ell_j}$ after the same sum; and the bracket is \Cref{rot:hyb:lem:CL} for the $C^1$ field $\partial_{y''}\tilde S^S_j(\cdot,y'')\,u$ (\Cref{rot:hyb:lem:ESB}(b),(c): $P_0\le C\ell_j^{-2}e^{-\delta\dist/\ell_j}$, $P_1\le C\ell_j^{-3}e^{-\delta\dist/\ell_j}$), $\le C(\eta'/\ell_j)\ell_j^{-1}e^{-\delta\dist/\ell_j}$. The weakest gain is the first sum's, $(\eta'/\ell_j)^{\theta_\partial}$. \emph{$j\in\{0,1\}$}: each line, also with the derivative, is bounded separately by $Ce^{-\delta\dist/\ell_j}$ (times $\ell_j^{-1}$ with the derivative) (\cite[(3.64)]{K1} for $Y$; \Cref{rot:hyb:lem:A0}(a) summed along $\le10.5\ell_j$ for $X$), and $(\eta'/\ell_j)^{1/2}\ge5^{-1/2}$: there is no gain and the factor $\eta'^\theta$ is paid by a degree loss, exactly as in \Cref{rot:cor:unif}(a).
\end{proof}

\begin{lemma}[double contour lines]\label{rot:hyb:lem:DCL}
For $0\le j<r$ and $x_1,x_2\in\Lambda_F$ let
\[
\cD^F_j(x_1;x_2)\defeq\sum_{b_1\in\Gamma^{F,(j+1)}_{x_{1,j+1},x_1}}\sum_{b_2\in\Gamma^{F,(j+1)}_{x_{2,j+1},x_2}}\eta'^2S^F_j(x_{b_1},x_{b_2})\,\hat b_1\cdot\hat b_2,
\]
the line joining two contour factors of \cite[(3.26)]{K1} (at one vertex with $m\ge2$, at the two factors of a product (3.26)$\times$(3.26), or at two vertices). Then $|\cD^F_j(x_1;x_2)|\le C\ell_je^{-\delta|x_{1,j+1}-x_{2,j+1}|/\ell_j}$ for both flows, and at matched points $|\cD^X_j(x_1';x_2')-\cD^Y_j(x_1'';x_2'')|\le C(\eta'/\ell_j)^{1/2}\ell_je^{-\delta|x_{1,j+1}-x_{2,j+1}|/\ell_j}$. In Ba{\l}aban's degree $D(H)=\sum_ld_l+\sum_va_v+3(|V(H)|-1)$ (\cite[(2.1)--(2.2)]{Ba3}, with the two summation points of a block product as vertices, as in \Cref{rot:cor:OF}) a double line has $d_{l_0}=+1$, and every connected all-scale-summed subgraph $H$ containing it has $D(H)\ge1$: the double line is never the binding constraint of the weakened-element bound.
\end{lemma}
\begin{proof}
$\cD^F_j(x_1;x_2)=\sum_{b_2}\eta'\hat b_2\cdot\cL^F_j(x_1;x_{b_2})$. The size bound: lengths $\le10.5\ell_j$ ($X$), $\le7.5\ell_j$ ($Y$) (\Cref{rot:hyb:lem:len}(i)), \Cref{rot:hyb:lem:A0}(a), \cite[(3.63)]{K1}; since both contours lie in open rotated cubes of side $5\ell_j$ about their roots, $e^{-\delta|x_{b_1}-x_{b_2}|/\ell_j}\le e^{5\sqrt3\delta}e^{-\delta|x_{1,j+1}-x_{2,j+1}|/\ell_j}$. For $j\ge2$ put $\cV(r)\defeq\sum_{b_1''}\eta'\,I_S[S^Y_j(x_{b_1''},\cdot)](r)\,\hat b_1''$ (extension in the second argument); then
\[
\cD^X-\cD^Y=\sum_{b_2}\eta'\hat b_2\cdot[\cL^X_j(x_1';x_{b_2})-\cL^Y_j(x_1'';\tau x_{b_2})]+\sum_{b_2}\eta'\hat b_2\cdot[\cV(\tau x_{b_2})-\cV(x_{b_2})]+[a^X(\Gamma^X_2)-a^Y(\Gamma^Y_2)]_{\cA=\cV}
\]
exactly (\Cref{rot:lem:vert}(i); $\tau x_{b_2}\in\Lambda_Y$); the three terms are bounded by \Cref{rot:hyb:lem:CLN} summed along $\Gamma^X_2$, by $c_\tau\eta'\sup|\nabla\cV|$ with $|\nabla\cV|\le C\ell_j^{-1}$ (\Cref{rot:hyb:lem:ESH}(iii), \Cref{rot:hyb:lem:ESB}(b)), and by \Cref{rot:hyb:lem:CL} for partial contours with $P_0\le C$, $P_1\le C\ell_j^{-1}$. For $j\le1$ use the size bounds and $(\eta'/\ell_j)^{1/2}\ge5^{-1/2}$. Degree: the size $C\ell_j$ against $\ell_j^{-1}$ for a slice line ($d_l=-1$) and $\ell_j^0$ for a one-ended contour line ($d_l=0$) gives $d_{l_0}=+1$ (\Cref{rot:cor:OF}: $d_l=-1+\#(\text{contour ends})$). For a connected all-scale-summed $H$ containing $l_0$: if $H\setminus l_0$ is connected (or $l_0$ is a self-loop), $D(H)=D(H\setminus l_0)+1\ge1$, because $D\ge0$ for every all-scale-summed subgraph that does not contain $l_0$ (the per-flow renormalised structure, \Cref{rot:cor:OF}); if $l_0$ is a bridge between $H_1$ and $H_2$, $D(H)=D(H_1)+D(H_2)+1+3\ge4$. So $D(H)-\theta\ge\frac12$ for $\theta=\frac12$.
\end{proof}

\begin{proposition}[every difference element at $A\ne0$]\label{rot:hyb:prop:list}
In the setting of \Cref{rot:prop:BD}(b), after \Cref{rot:lem:CB}, order the elements of a renormalised family consistently across the two flows and telescope as in \Cref{rot:prop:cp}. Every difference element $e^X_i-e^Y_i$ is of one of the following types, with the stated bound:
\begin{enumerate}[label=\textup{(\roman*)},leftmargin=*]
\item an $A=0$ line: a fine slice $S^F_j$, $j<r$ (\Cref{rot:cor:unif}(a): $C_D\eta'^{\theta}\ell_j^{-1-\theta}e^{-\delta|x'-y'|/\ell_j}$, $\theta=\frac12$), a coarse slice or $\Gamma^F_n$ (\Cref{rot:cor:unif}(b): $C\eta^{1/2}\ell_j^{-1}e^{-\delta|x'-y'|/\ell_j}$), the coarse multiplier (\Cref{rot:prop:fc}: $0.2a\eta^2$ relative), on $\Lambda_F$;
\item a contour line (3.27): \Cref{rot:hyb:lem:CLN}, gain $(\eta'/\ell_j)^{1/2}$; with one lattice derivative at its lattice end, gain $(\eta'/\ell_j)^{\theta_\partial}$, $\theta_\partial=\alpha/(2(2+\alpha))$ (\Cref{rot:hyb:lem:DS});
\item a contour vertex (3.24)--(3.26): \Cref{rot:hyb:lem:CV}, entrywise $\le Ce_rp_r^2\eta^{1/2}$, relative gain $\eta^{1/2}$;
\item an external lifted field $\tilde A^{(r),F}$, $H^F\phi_r$ or its covariant derivative: \Cref{rot:lem:H,rot:lem:grad}, $\le C\eta^{1/2}$ times the field norm;
\item a constant-background phase (3.46): \Cref{rot:hyb:lem:bph}, $\le e\,p_r\lmu^{-1}c_\tau\delta_r^{-1/2}\eta'$;
\item a hopping vertex of $V_r(A',0)$ or a vertex function (couplings, partition of unity, the Cameron--Martin form): \Cref{rot:lem:vert}(ii),(iii), unchanged;
\item a double contour line, joining two contour factors of (3.26): \Cref{rot:hyb:lem:DCL}, size $\le C\ell_j$, relative gain $(\eta'/\ell_j)^{1/2}$ on slice $j$.
\end{enumerate}
Each weakened element therefore obeys the bound of its type weakened by the degree $\theta=\frac12$ (by $\theta_\partial\ge\theta_1$ for a contour line carrying a lattice derivative, which is all \Cref{rot:thm:H1S} asks) with a constant $\le C\,p_r^{2}\,\delta_r^{-1/2}$ (the worst of (iii) and (v)); the per-flow remainders (\Cref{rot:lem:CB}(2)) contribute $(Ce_rp_r^2)^{n+1}+e^{-\delta'p_r}$ per graph.
\end{proposition}
\begin{proof}
\Cref{rot:lem:CB} lists what carries the background after the reduction; the compared lines are $A=0$ lines, the contour lines, one-ended (3.27) and double (both ends on contours), are the only lines whose shape depends on the flow, and the only background-dependent vertices are the contour vertices (3.24)--(3.26) (\cite[p.~658]{K1}: ``The remaining vertices occur when the renormalization transformation of the scalar field is expanded about $A'=0$'') and the phases (3.46). The external coarse fields $A_r(w_i)$ of (3.55) are common. The bounds are the cited lemmas.
\end{proof}

\begin{proof}[Proof of \Cref{rot:prop:BD}(b)]
The difference elements $\Delta\Gamma$, $\Delta H$, $\Delta(A^{-1})$ at regular $A$ are, after the reduction, the $A=0$ differences of \Cref{rot:cor:unif} (types (i) and (iv) of \Cref{rot:hyb:prop:list}), with the factor $\eta'^{1/2}$. The contour vertices at the reduced background differ entrywise by a relative $\eta^{1/2}$ against their size $Ce_rp_r^2$ (\Cref{rot:hyb:lem:CV}); the constant-background phases at matched vertices differ by $\le e\,p_r\lmu^{-1}c_\tau\delta_r^{-1/2}\eta'$ (\Cref{rot:hyb:lem:bph}); the per-slice pieces carried by the contour vertex, with one end on a contour (\cite[(3.27)]{K1}) or both ends on contours, obey the slice-$j$ bounds \cite[(3.63)--(3.65)]{K1} in both flows ((a) for the hybrid flow, \cite[Prop.~3.7]{K1} for the standard flow) with the relative mismatch $(\eta'/\ell_j)^{1/2}$ (\Cref{rot:hyb:lem:CLN,rot:hyb:lem:DCL}), the one-ended lines also with one lattice derivative at the lattice end, there with the relative mismatch $(\eta'/\ell_j)^{\theta_\partial}$ of \Cref{rot:hyb:lem:DS} (\Cref{rot:hyb:lem:CLN}; a double line has no lattice end). A constant background is pure gauge on the infinite lattice and gives the same phase, up to the matching segment, for every contour (\Cref{rot:hyb:lem:bph}); what would break (b) are mis-paired contours or mis-rooted blocks, both excluded by \Cref{rot:def:flows} and \Cref{rot:conv:C}(i). The gains are not uniform in the size of the constant part of the background: they carry the small-field bound \cite[(3.2)]{K1}, as King's own reduction does.
\end{proof}
